\chapter{Suy luận thống kê: Odds Ratio (OR)}

\section{Từ hệ số \texorpdfstring{$\beta$}{beta} đến Odds Ratio}
Trong mô hình logit:
\[
\log\left(\frac{p}{1-p}\right) = \beta_0 + \sum_{j=1}^{d}\beta_j x_j.
\]
Với mỗi biến \(x_j\), \textit{Odds Ratio} (tỷ số odds) được định nghĩa:
\[
OR_j = \exp(\beta_j).
\]
Diễn giải: nếu tăng \(x_j\) thêm 1 đơn vị (giữ các biến khác không đổi), odds nghỉ việc sẽ nhân với \(OR_j\).

\section{Kiểm định giả thuyết cho từng hệ số}
Khi báo cáo p-value cho từng hệ số \(\beta_j\), ta thực chất đang kiểm định:
\[
H_0: \beta_j = 0 \quad \Leftrightarrow \quad OR_j = 1
\]
chống lại giả thuyết đối:
\[
H_1: \beta_j \neq 0 \quad \Leftrightarrow \quad OR_j \neq 1.
\]
Với mức ý nghĩa \(\alpha = 0.05\), nếu p-value \(< \alpha\) thì bác bỏ \(H_0\), kết luận biến \(x_j\) có ý nghĩa thống kê (tác động đáng kể lên log-odds nghỉ việc); ngược lại, không đủ bằng chứng để khẳng định biến đó khác 0 trong mô hình hiện tại.

\section{Baseline (nhóm tham chiếu) cho biến phân loại}
Do categorical được dummy hoá theo cơ chế \textit{drop-first} (loại một nhóm làm baseline), các biến:
\begin{itemize}
  \item \texttt{OverTime\_Yes} được hiểu \textbf{so với baseline OverTime\_No}.
  \item \texttt{Gender\_Male} được hiểu \textbf{so với baseline Gender\_Female}.
  \item Các JobRole (như \texttt{JobRole\_Sales Representative}) so với \textbf{baseline JobRole}.
\end{itemize}

\section{Bảng OR và Forest plot}
\begin{longtable}{@{}p{4.2cm}p{2.4cm}p{5.2cm}p{2.2cm}@{}}
  \caption{Bảng Odds Ratio (OR), 95\% CI và p-value} \\
  \toprule
  \textbf{Feature} & \textbf{OR} & \textbf{95\% CI} & \textbf{p-value} \\
  \midrule
  \endfirsthead
  \toprule
  \textbf{Feature} & \textbf{OR} & \textbf{95\% CI} & \textbf{p-value} \\
  \midrule
  \endhead
  OverTime\_Yes & 7.2785 & [4.9777,\;10.6429] & \(1.33\times 10^{-24}\) \\
  EnvironmentSatisfaction & 0.6476 & [0.5506,\;0.7618] & \(1.56\times 10^{-7}\) \\
  NumCompaniesWorked & 1.2174 & [1.1284,\;1.3134] & \(3.81\times 10^{-7}\) \\
  JobSatisfaction & 0.6614 & [0.5638,\;0.7759] & \(3.84\times 10^{-7}\) \\
  BusinessTravel\_Travel\_Frequently & 6.7167 & [2.9945,\;15.0656] & \(3.82\times 10^{-6}\) \\
  JobInvolvement & 0.5892 & [0.4635,\;0.7488] & \(1.53\times 10^{-5}\) \\
  DistanceFromHome & 1.0465 & [1.0246,\;1.0688] & \(2.44\times 10^{-5}\) \\
  YearsSinceLastPromotion & 1.1915 & [1.0969,\;1.2943] & \(3.31\times 10^{-5}\) \\
  JobRole\_Sales Representative & 7.1923 & [2.4411,\;21.1910] & \(3.45\times 10^{-4}\) \\
  MaritalStatus\_Single & 3.1158 & [1.5833,\;6.1315] & 0.0010 \\
  YearsInCurrentRole & 0.8644 & [0.7910,\;0.9446] & 0.0013 \\
  RelationshipSatisfaction & 0.7748 & [0.6593,\;0.9105] & 0.0019 \\
  JobRole\_Laboratory Technician & 4.3769 & [1.6963,\;11.2937] & 0.0023 \\
  WorkLifeBalance & 0.6895 & [0.5408,\;0.8793] & 0.0027 \\
  YearsWithCurrManager & 0.8763 & [0.7991,\;0.9610] & 0.0050 \\
  BusinessTravel\_Travel\_Rarely & 2.7766 & [1.3189,\;5.8454] & 0.0072 \\
  TrainingTimesLastYear & 0.8295 & [0.7190,\;0.9570] & 0.0104 \\
  YearsAtCompany & 1.0944 & [1.0148,\;1.1803] & 0.0192 \\
  JobRole\_Sales Executive & 2.7824 & [1.1577,\;6.6869] & 0.0222 \\
  Age & 0.9695 & [0.9441,\;0.9956] & 0.0222 \\
  Gender\_Male & 1.4830 & [1.0328,\;2.1295] & 0.0328 \\
  TotalWorkingYears & 0.9409 & [0.8885,\;0.9964] & 0.0371 \\
  JobRole\_Human Resources & 3.7587 & [0.9984,\;14.1499] & 0.0503 \\
  StockOptionLevel & 0.8069 & [0.5926,\;1.0986] & 0.1730 \\
  DailyRate & 0.9997169 & [0.9992858,\;1.0001483] & 0.1983 \\
  const & 7.4013 & [0.3229,\;169.6507] & 0.2104 \\
  MaritalStatus\_Married & 1.3812 & [0.8185,\;2.3308] & 0.2264 \\
  JobRole\_Research Director & 0.3420 & [0.0477,\;2.4521] & 0.2857 \\
  JobRole\_Research Scientist & 1.6960 & [0.6427,\;4.4752] & 0.2860 \\
  EducationField\_Medical & 0.5117 & [0.1162,\;2.2529] & 0.3756 \\
  EducationField\_Other & 0.5203 & [0.1036,\;2.6134] & 0.4276 \\
  EducationField\_Life Sciences & 0.5638 & [0.1280,\;2.4821] & 0.4485 \\
  PercentSalaryHike & 0.9783 & [0.9059,\;1.0564] & 0.5755 \\
  EducationField\_Technical Degree & 1.4123 & [0.3072,\;6.4929] & 0.6574 \\
  JobRole\_Manufacturing Director & 1.2463 & [0.4387,\;3.5405] & 0.6794 \\
  MonthlyRate & 1.0000051 & [0.9999806,\;1.0000296] & 0.6833 \\
  PerformanceRating & 1.1346 & [0.5197,\;2.4772] & 0.7512 \\
  JobLevel & 0.9135 & [0.4914,\;1.6981] & 0.7748 \\
  HourlyRate & 1.0011 & [0.9925,\;1.0098] & 0.8085 \\
  EducationField\_Marketing & 0.8415 & [0.1737,\;4.0775] & 0.8303 \\
  JobRole\_Manager & 1.1686 & [0.2495,\;5.4729] & 0.8433 \\
  MonthlyIncome & 1.0000111 & [0.9998515,\;1.0001708] & 0.8912 \\
  Education & 1.0025 & [0.8442,\;1.1905] & 0.9769 \\
  \bottomrule
\end{longtable}


\begin{figure}[H]
  \centering
  \includegraphics[width=0.9\textwidth]{forest_plot_odds_ratio.png}
  \caption{Forest plot Odds Ratio}
\end{figure}

\noindent
Trong bảng OR, mỗi dòng tương ứng với một biến (hoặc nhóm tham chiếu dummy) trong mô hình logit.
Forest plot trực quan hoá cùng một thông tin: trục hoành ở thang log hoặc thang OR, với điểm giữa là OR ước lượng, đoạn thẳng ngang biểu diễn 95\% CI và đường dọc tại OR = 1 là ranh giới ``không hiệu ứng''. Các biến có đoạn CI không cắt qua 1 và nằm xa về một phía thể hiện tác động mạnh và ổn định lên quyết định nghỉ việc.

\section{Một số điểm diễn giải trọng tâm}
\begin{itemize}
  \item \textbf{OverTime\_Yes}: OR cực kỳ cao (quanh mức 7.2). Nghĩa là nhân viên phải làm thêm giờ có odds nghỉ việc cao gấp $\approx$ 7 lần so với người không làm thêm giờ (với các điều kiện khác không đổi).
  \item \textbf{JobRole\_Sales Representative}: OR cao (quanh 7.1), cho thấy rủi ro nghỉ việc đặc thù rất lớn ở vị trí nhân viên bán hàng.
  \item \textbf{EnvironmentSatisfaction}: OR $\approx$ 0.64. Tăng 1 điểm hài lòng về môi trường làm việc làm giảm $36\%$ odds nghỉ việc.
  \item \textbf{JobInvolvement}: OR $\approx$ 0.58. Sự gắn kết với công việc tăng lên giúp giữ chân nhân viên hiệu quả.
\end{itemize}

\noindent
Nhìn tổng thể bảng OR, có thể nhóm các biến thành hai loại chính: (i) \textit{yếu tố rủi ro} với OR $>$ 1 làm tăng odds nghỉ việc (OverTime, Sales Representative, DistanceFromHome), và (ii) \textit{yếu tố bảo vệ} với OR $<$ 1 làm giảm odds nghỉ việc (EnvironmentSatisfaction, JobInvolvement, RelationshipSatisfaction). Các biến liên tục như \texttt{MonthlyIncome} có OR rất gần 1 vì đơn vị tính là 1 USD; trên thực tế, tác động biên của 1000 USD sẽ lớn hơn rất nhiều.

Các yếu tố như OverTime và EnvironmentSatisfaction đều có p-value cực nhỏ, cho thấy ảnh hưởng của chúng lên quyết định nghỉ việc là vô cùng ổn định. Từ góc độ nghiệp vụ (HR), \textit{OverTime} là chỉ báo rủi ro cần được giám sát nghiêm ngặt, trong khi việc nâng cao \textit{EnvironmentSatisfaction} là đòn bẩy chính sách lâu dài.

\section{Phương trình logit ước lượng}
Về mặt lý thuyết, mô hình logit giả định rằng log-odds nghỉ việc là một hàm tuyến tính lập trên ma trận thiết kế đã dummy hoá:
\[
\log\left(\frac{p}{1-p}\right) = \beta_0 + \sum_{j=1}^{d}\beta_j x_j.
\]
Các hệ số $\hat{\beta}_j$ chính là $\ln(OR_j)$. Vì số lượng biến quá lớn, việc trình bày toàn bộ phương trình sẽ gây loãng thông tin. Bảng Odds Ratio phía trên và \texttt{assets/outputs/logit\_summary.txt} cung cấp đầy đủ bức tranh định lượng để tra cứu hệ số chính xác. Một nhân viên có $x_j$ nhất định sẽ có xác suất rời công ty được tính qua:
\[
p = \sigma(\text{logit}) = \frac{1}{1 + e^{-\text{logit}}}.
\]
